\section{The cross-generality objection --- and why $\Prin$ answers it}
\label{sec:cross}

\textbf{Objection.} Some skills \emph{within} a domain --- experiment design, tool orchestration, recognizing when the data violate the domain's rules --- appear to require cross-domain generality. If so, \DCI\ is not a clean stepping stone but a partial one.

\textbf{Resolution.} The objection conflates two axes that a vertically organizing principle separates. Verticality answers each on different grounds.

\emph{Epistemic axis (cross-scale transfer and violation detection).} Because $\Prin$ is the scale-invariant of $\Dom$, a rule-violation \emph{is} a measurable departure from $\Prin$: an invariant is, by construction, the object whose violation is diagnostic. No external reasoner is required to flag it. Likewise, transfer across scales is not bolted on --- it is constitutive of holding $\Prin$ at all, since $\Prin$ has the same form at every level. The effective scope of a $\Prin$-grounded \DCI\ is thus the \emph{entire vertical hierarchy} $\Dom$ spans, not a thin slice of it. Most of the ``isn't this just narrow AI'' reflex dissolves here.

\emph{Pragmatic/meta axis (search, experiment design, orchestration).} These are not cross-scale but meta-level, and scale-invariance alone does not cover them. They are covered by a further, substantive property: \textbf{$\Prin$ is reflexive across the object/meta boundary} --- it governs not only the dynamics of $\Dom$ but the dynamics of \emph{discovery over} $\Dom$. If skill acquisition is itself an action-minimizing trajectory, then experiment design is ``select the path through hypothesis space that minimizes network-weighted action,'' an instance of $\Prin$ at the meta-level rather than a separate faculty demanding external generality.

This reflexivity is not asserted; it is demonstrated, and along two independent routes (Figure~\ref{fig:reflexivity}) --- with the neural-architecture paper \citep{frasch2026c} supplying the engineering proof of the first and the formal basis of the second.

\emph{Instance 1 --- reflexivity through search.} The \textbf{Emin/Imax} formalism maps \NWAP\ onto neural architecture design: an energy-first objective (the action functional) evaluated across \textbf{2{,}203 experiments} spanning vision, text, neuromorphic, and physiological data, with $10$ seeds per configuration and a factorial analysis \citep{frasch2026c}. The \emph{same} principle that describes an object system also governs the \emph{search} for architectures that model it --- a strong architecture$\times$dataset interaction ($\etap = 0.44$) shows the objective selecting task-appropriate structure rather than a fixed winner. Here the meta-level is design over model space, and $\Prin$ governs it.

\emph{Instance 2 --- reflexivity through inference.} \citet{frasch2026c} also places \NWAP\ in the same variational equivalence class as the free-energy principle and the evidence lower bound (ELBO) of Bayesian inference --- three trade-off-based objectives that balance a data-fit term against a coupling/complexity cost. The paper is careful here, and so is this argument: \NWAP\ is \emph{not} a probabilistic ELBO, but a member of the same equivalence class of variational principles. That class-membership is what matters. The ELBO and the free-energy principle are not object-level descriptions; they \emph{are} the operations by which a system infers a generative model of its data. If $\Prin$ shares their variational form, then model-discovery over the domain is governed by an objective of the \emph{same class} as $\Prin$ --- structural kinship at the object/meta boundary, not a loose metaphor and not a claim of literal identity. And this is not only formal: the same action-functional constraint has been validated at biology scale, predicting the emergent modular organization of real marine metabolic networks, with a modularity excess of $\dQ \approx 0.15$--$0.40$ over configuration-model, label-permutation, and bipartite-incidence null models \citep{frasch2026d}. So reflexivity is witnessed both in \emph{search} (in silico, NAS) and in \emph{inference} (empirically, on real data) --- two different meta-level operations, two different evidentiary bases, with \citeyear{frasch2026c} as the shared formal spine.

\begin{figure*}[t]
\centering
\begin{tikzpicture}[font=\small,>={Latex[length=2.4mm]},
  box/.style={draw,rounded corners,align=center,inner sep=5pt,minimum height=11mm},
  lvl/.style={box,fill=NavyBlue!7,text width=4.5cm},
  ver/.style={box,fill=gray!12,text width=3.5cm},
  prin/.style={draw,circle,fill=orange!22,inner sep=2pt,minimum size=12mm,font=\normalsize}]
\node[prin] (pi) {$\Prin$};
\node[lvl,above right=4mm and 22mm of pi] (obj) {\textbf{Object level}\\ domain $\Dom$: dynamics extremize $\Prin$};
\node[lvl,below right=4mm and 22mm of pi] (meta) {\textbf{Meta level}: discovery over $\Dom$\\ search $+$ inference; objective $\in$ class$(\Prin)$};
\node[ver,right=24mm of meta] (ver) {\textbf{Verifier} $\Rule$\\ nomological machine};
% reflexivity: one principle governs both levels
\draw[->] (pi.north) to[out=90,in=180] node[above,font=\footnotesize]{describes dynamics} (obj.west);
\draw[->] (pi.south) to[out=-90,in=180] node[below,font=\footnotesize]{governs discovery} (meta.west);
% reflexive object/meta coupling
\draw[->] (meta.north) -- node[right,font=\footnotesize]{proposes law/model} (obj.south);
% closed discovery loop against the verifier
\draw[->] (meta.east) -- node[above,font=\footnotesize]{candidate} (ver.west);
\draw[->] (ver.south) to[out=-90,in=-90] node[below,font=\footnotesize]{adjudicates} (meta.south);
\draw[->,dashed] (obj.east) to[out=0,in=90] node[above right,font=\footnotesize]{ground truth} (ver.north);
\end{tikzpicture}
\caption{The object/meta reflexivity loop. A single vertically organizing principle $\Prin$ both \emph{describes} the object-level dynamics of $\Dom$ and \emph{governs} the meta-level discovery over $\Dom$ --- the reflexivity that answers the cross-generality objection. Discovery is witnessed along two routes: Instance~1, search (energy-first NAS, \citealp{frasch2026c}), and Instance~2, inference (\NWAP\ in the free-energy/ELBO variational class, \citealp{frasch2026c}, validated on real data in \citealp{frasch2026d}). The executable verifier $\Rule$ --- a nomological machine (\S\ref{sec:lineage}) --- closes the loop by adjudicating candidates against ground truth, covering any residual the principle does not.}
\label{fig:reflexivity}
\end{figure*}

\textbf{Optimization is not yet agency.} We are careful not to overread the reflexivity evidence, because the object/meta step is where this argument is most easily overstated. Energy-minimization selecting a task-appropriate architecture (Instance~1) is advanced \emph{optimization} over a fixed design space; it is not, by itself, epistemic agency, and we do not claim the NAS result demonstrates the latter. Two meta-skills lie strictly beyond it: (a)~recognizing that the domain's own generative assumptions are wrong --- that it is the rule substrate, not the model, that needs revision --- and (b)~\emph{basis expansion}: inventing a genuinely new term when the existing candidate library cannot express the target law (the open step already flagged in \S\ref{sec:verifier}). What reflexivity actually buys is narrower and still substantive: the \emph{same} objective that scores object-level dynamics also governs search and inference \emph{within a fixed representational basis}. The move from selection to genuine discovery --- basis expansion --- is the sharpest open problem this frame names, not one it claims to have solved. In the interim the executable verifier is what carries discovery across that gap: it makes even an ill-motivated candidate term cheap to falsify, so the mechanism for extending the basis can be dumb (enumeration, mutation, LLM proposal) as long as $\Rule$ adjudicates it. A formal mechanism by which $\Prin$ itself \emph{generates} the missing term remains future work.

\textbf{Composition with the verifier.} Where even these leave residue --- messier real-world meta-skills such as physical experiment orchestration --- the verifier composes with $\Prin$: the executable substrate $\Rule$ substitutes for meta-judgment in exactly the cases where $\Prin$'s meta-coverage is thinnest. Reflexivity of $\Prin$ (search \emph{and} inference) \emph{plus} the verifier substrate together close the loop --- a more robust guarantee than any one of them alone.

The correct statement of the claim is therefore precise, not hand-waving: \emph{a vertically organizing principle addresses the cross-generality concern because it is reflexive across the object/meta boundary; the NAS \citep{frasch2026c} and biology-scale validation \citep{frasch2026d} results are independent evidence that it is; and the verifier substrate covers the residual.}
